\documentclass[11pt,a4paper]{article}
\usepackage[utf8]{inputenc}
\usepackage{indentfirst}
\usepackage{amsmath,amssymb}
\usepackage{graphicx}
\usepackage{listings}
\usepackage{xcolor}
\usepackage[colorlinks=true,linkcolor=blue]{hyperref}

\definecolor{lightgray}{rgb}{0.95,0.95,0.95}
\lstset{
    basicstyle=\ttfamily\small,
    backgroundcolor=\color{lightgray},
    breaklines=true,
    frame=single,
    framerule=0.5pt,
    numbers=left,
    numberstyle=\tiny,
    stepnumber=1,
    captionpos=t,
}

\appendix
\section{核心算法实现}

\subsection{模块一：五元组稳态检测}

\begin{lstlisting}[language=Python, caption={五元组稳态判据与Rolling MAD去噪核心}, label={lst:preprocessor}]
# 五元组物理阈值（五值全 > 0 表示正常推进）
STEADY_STATE_THRESHOLDS = {
    "总推进力": 0.0, "刀盘扭矩": 0.0, "推进速度": 0.0,
    "刀盘转速": 0.0, "贯入度": 0.0,
}

def _extract_steady_state_mask(chunk):
    """五值全 > 0 → 正常推进；任一为 0 → 停机"""
    return ((chunk["总推进力"] > 0) & (chunk["刀盘扭矩"] > 0) &
            (chunk["推进速度"] > 0) & (chunk["刀盘转速"] > 0) &
            (chunk["贯入度"] > 0))

def _rolling_mad_denoise(series, window=60, tau=3.5, k=1.4826):
    """Rolling MAD：局部中位数 ± 3.5σ 动态阈值去毛刺"""
    rolling_med = series.rolling(window, min_periods=1, center=True).median()
    abs_dev = (series - rolling_med).abs()
    rolling_mad = abs_dev.rolling(window, min_periods=1, center=True).median()
    robust_sigma = k * rolling_mad
    glitch_mask = abs_dev > tau * robust_sigma
    return series.ffill().bfill()  # 前向/后向填充，不删行
\end{lstlisting}

\subsection{模块二：7维物理不变量特征工程}

\begin{lstlisting}[language=Python, caption={物理导向7维特征衍生核心公式}, label={lst:feature_engineer}]
CUTTERHEAD_AREA = 188.7  # 刀盘截面积 (m²)
FINAL_7D = ['SE', 'FPI', 'TPI', 'P2.1泵电流', '主油箱油温', '泥水仓顶部1压力', '排浆密度']

def compute_physics_invariants(df):
    """注入三大高阶不变量（物理导向特征平替原则）"""
    Fv, T, v, n = df['总推进力'].values, df['刀盘扭矩'].values, df['推进速度'].values, df['刀盘转速'].values
    p = np.where(n > 0, v / n, np.nan)  # 贯入度 p = v / n

    # SE = Fv/A + 2πT/(A·p)    切削比能 [kPa]
    df['SE']  = (Fv / CUTTERHEAD_AREA) + (2 * np.pi * T) / (CUTTERHEAD_AREA * p)

    # FPI = Fv / p              推力贯入度指数 [kN/mm]
    df['FPI'] = Fv / p

    # TPI = T  / p              扭矩贯入度指数 [kN·m/mm]
    df['TPI'] = T / p

    return df[FINAL_7D]  # 严格剔除底层机械参数
\end{lstlisting}

\subsection{模块三：TS-VAE网络架构与ELBO损失}

\begin{lstlisting}[language=Python, caption={TS-VAE编码器-潜空间-解码器架构}, label={lst:tsvae}]
class TSVAE(nn.Module):
    """时序变分自编码器：输入 60×7 → 潜空间 16 → 重构 60×7"""
    def __init__(self, input_size=7, hidden=64, latent=16, seq_len=60, beta=0.001):
        super().__init__()
        self.encoder = nn.GRU(input_size, hidden, num_layers=2, batch_first=True)
        self.fc_mu    = nn.Linear(hidden, latent)
        self.fc_logvar = nn.Linear(hidden, latent)
        self.decoder = nn.GRU(latent, hidden, num_layers=2, batch_first=True)
        self.fc_out  = nn.Linear(hidden, input_size)
        self.latent_dim, self.seq_len, self.beta = latent, seq_len, beta

    def reparameterize(self, mu, logvar):
        return mu + torch.exp(0.5 * logvar) * torch.randn_like(mu)

    def forward(self, x):
        _, h = self.encoder(x)
        mu, logvar = self.fc_mu(h[-1]), self.fc_logvar(h[-1])
        z = self.reparameterize(mu, logvar)
        h0 = torch.tanh(nn.Linear(self.latent_dim, 64*2)(z)).view(2, -1, 64)
        out = self.decoder(z.unsqueeze(1).expand(-1, self.seq_len, -1), h0)
        return self.fc_out(out), mu, logvar, z

    def elbo_loss(self, x, recon, mu, logvar):
        """ELBO = L_recon + β·L_KL，β=0.001 弱化KL防止后验崩塌"""
        recon_loss = F.mse_loss(recon, x, reduction='sum') / x.size(0)
        kl_loss = -0.5 * torch.sum(1 + logvar - mu.pow(2) - logvar.exp()) / x.size(0)
        return recon_loss + self.beta * kl_loss, recon_loss, kl_loss
\end{lstlisting}

\subsection{物理场分组与相关性配置}

\begin{lstlisting}[language=Python, caption={跨物理场MSE分组与Spearman相关矩阵}, label={lst:phy_groups}]
# 四大物理场分组（Fig 11热力图）
PHY_GROUPS = {
    '能效/机械': ['SE', 'FPI', 'TPI'],
    '电气': ['P2.1泵电流'],
    '热力': ['主油箱油温'],
    '流体': ['泥水仓顶部1压力', '排浆密度'],
}

# 7D Spearman相关矩阵（设计层面约束妥协：SE↔TPI=0.941）
#              SE     FPI    TPI    I_p21  T_oil  P_top  rho
# SE          1.000  0.765  0.941  0.029 -0.120  0.006  0.028
# FPI         0.765  1.000  0.774 -0.191 -0.160 -0.010 -0.211
# TPI         0.941  0.774  1.000  0.010  0.030 -0.015  0.011
# I_p21       0.029 -0.191  0.010  1.000  0.085 -0.053  0.504
# T_oil      -0.120 -0.160  0.030  0.085  1.000 -0.000  0.151
# P_top       0.006 -0.010 -0.015 -0.053 -0.000  1.000  0.007
# rho_out     0.028 -0.211  0.011  0.504  0.151  0.007  1.000
# 跨域特征对相关性均 < 0.2，满足独立性要求
\end{lstlisting}

\end{document}